\clearpage\thispagestyle{otFront}\phantomsection\label{parcours}\pdfbookmark[0]{Construire son parcours}{bm-parcours}\otPageTitre{Construire son parcours}{Une année avec des retours réguliers}\fontsize{11}{15}\selectfont

\begin{tabularx}{\linewidth}{@{}lY@{}}\toprule
Calculer & 1 à 7 : opérations, divisibilité, relatifs, fractions, carrés et cubes.\\[3mm]
Modéliser & 8 à 12 : formules, distributivité, équations simples, proportionnalité, repérage et dépendance.\\[3mm]
Construire & 13 à 17 : angles, triangles, symétries, parallélogrammes, solides et volumes.\\[3mm]
Interpréter & 18 et 19 : données, moyenne simple et probabilités.\\[3mm]
Programmer & 20 : séquences, expressions et répétitions, à reprendre pendant l’année.\\\bottomrule
\end{tabularx}
\begin{savoirfaire}[Une séance régulière]Relire un exemple, chercher deux
exercices sans correction, puis comparer les démarches. Refaire une
question manquée et noter la notion à revoir.\end{savoirfaire}
\begin{savoirfaire}[Revenir et mélanger]Reprendre les notions quelques jours
puis quelques semaines après. Mélanger calcul, graphique et justification
pour apprendre à choisir un outil sans dépendre du titre du chapitre.\end{savoirfaire}
Les défis demandent environ 30 à 60 minutes, éventuellement en plusieurs
séances. Les bilans sont prévus pour deux heures. Adapter le temps à
l’objectif de travail ; le tableau de bord sert à suivre les reprises.

\clearpage
